\chapter{Kongruenzen der Verbandsoperationen}
\label{b45guide:ch:congruences}

\section{Verträgliche Identifizierung von Elementen}

\begin{definition}[Halbverbands- und Verbandskongruenz]
\label{b45guide:def:congruence}
Eine \emph{Kongruenz} eines Halbverbandes $(S,\vee)$ ist eine
Äquivalenzrelation $\theta$ auf $S$, für die
\[
 x\mathrel\theta x',\ y\mathrel\theta y'
 \quad\Longrightarrow\quad
 x\vee y\mathrel\theta x'\vee y'
\]
gilt. Eine \emph{Verbandskongruenz} erfüllt zusätzlich dieselbe
Verträglichkeit für $\wedge$.
Die Klasse von $x$ wird mit $[x]_\theta$ bezeichnet.
Für einen Homomorphismus $f$ ist
\[
 \ker f=\{(x,y)\mid f(x)=f(y)\}
\]
seine \emph{Kernkongruenz}.
\end{definition}

\begin{theorem}[Quotienten und Homomorphiesatz]
\label{b45guide:thm:quotient}
Ist $\theta$ eine Halbverbandskongruenz, so ist
\[
 [x]_\theta\vee[y]_\theta=[x\vee y]_\theta
\]
eine wohldefinierte Halbverbandsoperation auf $S/\theta$.
Ist $\theta$ eine Verbandskongruenz, so liefert zusätzlich
\[
 [x]_\theta\wedge[y]_\theta=[x\wedge y]_\theta
\]
eine Verbandsstruktur. Die natürliche Abbildung $x\mapsto[x]_\theta$
erhält die Operationen.

Der Kern jedes Homomorphismus ist eine Kongruenz. Für jeden
Halbverbands- oder Verbandshomomorphismus $f:S\to T$ induziert
\[
 S/\ker f\longrightarrow f[S],\qquad [x]_{\ker f}\longmapsto f(x)
\]
einen Isomorphismus der jeweiligen Strukturen.
Mitgeführte Konstanten werden auf ihre Klassen abgebildet. \cite{nation-con}
\end{theorem}

\begin{theorem}[Prüfung durch Translationen]
\label{b45guide:thm:translations}
Eine Äquivalenzrelation $\theta$ auf einem Halbverband ist genau
dann eine Halbverbandskongruenz, wenn für jedes $c$ gilt
\[
 x\mathrel\theta y\quad\Longrightarrow\quad
 x\vee c\mathrel\theta y\vee c.
\]
Auf einem Verband müssen zusätzlich alle Abbildungen
$x\mapsto x\wedge c$ die Relation erhalten.
Damit genügt es, Verträglichkeit mit je einem festgehaltenen
Argument zu prüfen.
\end{theorem}

\section{Besonderheiten von Verbandskongruenzen}

\begin{theorem}[Klassen sind konvexe Unterverbände]
\label{b45guide:thm:classes}
Jede Kongruenzklasse eines Verbandes ist ein nichtleerer Unterverband.
Sie ist \emph{ordnungskonvex}, das heißt
\[
 a\leq z\leq b,\quad a\mathrel\theta b
 \quad\Longrightarrow\quad z\mathrel\theta a.
\]
In einem endlichen Verband ist daher jede Klasse ein Intervall
$[u,v]_L$. Darüber hinaus gilt in jedem Verband
\[
 x\mathrel\theta y
 \quad\Longleftrightarrow\quad
 x\wedge y\mathrel\theta x\vee y .
\]
Eine bloße Zerlegung in konvexe Unterverbände reicht im Allgemeinen
nicht aus, um eine Verbandskongruenz zu erhalten. \cite{nation-con}
\end{theorem}

\begin{definition}[Kongruenzverband]
\label{b45guide:def:con}
$\operatorname{Con}L$ bezeichnet die Menge aller Verbandskongruenzen
von $L$. Für $\alpha,\beta\in\operatorname{Con}L$ setzen wir
\[
 \alpha\wedge_c\beta=\alpha\cap\beta.
\]
$\alpha\vee_c\beta$ ist die kleinste Verbandskongruenz, die
$\alpha\cup\beta$ enthält; sie ist der Durchschnitt aller
Verbandskongruenzen mit dieser Eigenschaft.
Die Indizes unterscheiden diese Operationen von $\wedge,\vee$ in $L$.
\end{definition}

\begin{theorem}[Funayama--Nakayama]
\label{b45guide:thm:funayama}
Für jeden Verband $L$ bilden diese Operationen auf
$\operatorname{Con}L$ einen Verband und erfüllen
\[
\begin{aligned}
 \alpha\wedge_c(\beta\vee_c\gamma)
   &=(\alpha\wedge_c\beta)\vee_c(\alpha\wedge_c\gamma),\\
 \alpha\vee_c(\beta\wedge_c\gamma)
   &=(\alpha\vee_c\beta)\wedge_c(\alpha\vee_c\gamma)
\end{aligned}
\]
für alle $\alpha,\beta,\gamma\in\operatorname{Con}L$.
Diese Distributivitätsgesetze sind eine Folgerung für die
\emph{Kongruenzen jedes Verbandes}; Distributivität von $L$
wird nicht vorausgesetzt. Die entsprechende Behauptung gilt
für Kongruenzverbände beliebiger Halbverbände nicht. \cite{nation-con,freese-nation}
\end{theorem}

\begin{theorem}[Endliche Verbandskongruenzen durch kollabierte Kanten]
\label{b45guide:thm:edges}
Sei $L$ endlich. Für $a<b$ schreiben wir $a\prec b$, wenn es
kein $z$ mit $a<z<b$ gibt.
Zwei Kongruenzen $\alpha,\beta$ von $L$ sind genau dann gleich,
wenn sie dieselben solchen Paare identifizieren:
\[
 \forall a\prec b\quad
 \bigl(a\mathrel\alpha b\iff a\mathrel\beta b\bigr).
\]
Diese Aussage liefert ein Eindeutigkeitskriterium;
nicht jede beliebige Auswahl von Kanten definiert eine Kongruenz.
\end{theorem}
